\section{Bounded cone and $F_4$ strata algorithms}
\label{app:cone-model-algorithms}

This appendix gives the exhaustive certification.  Root data are generated
from the $F_4$ Cartan matrix.  The reviewed fine-pair list is checked directly
and reconstructed when absent.  Roots are coefficient tuples in Bourbaki
order, and all rank calculations are exact.

\subsection{Field-degree bound}

For each of the $105$ downward-closed complements $K$, the program builds the
complete grouped cone matrix for the largest possible target set.  An integral
nonzero minor certifies its rank.  The resulting dimension estimate is affine
in $d=[F:\Q_p]$ and is automatic for $d\ge2$.  The contents of all witness
minors have prime support contained in $\{2,3\}$, so the estimate is unchanged
in the standing range $p>12$.

\begin{tcolorbox}
\begin{algorithmic}[1]
\Require The $F_4$ Cartan matrix.
\State Generate $\Phi^+$, its root order, and all $105$ complements $K$.
\For{each $K$}
  \State Build the full grouped cone matrix at $\Psi_2=\Phi^+$.
  \State Find an integral maximal nonzero minor and its bad primes.
  \State Solve the resulting affine dimension inequality for $d$.
\EndFor
\State Return the largest exceptional degree range.
\end{algorithmic}
\end{tcolorbox}
\renewcommand{\figurename}{Algorithm}
\captionof{figure}{A priori field-degree bound}
\label{alg:F4-degree-bound}

Thus only $d=1$ requires enumeration.

\subsection{Fine pairs}

At degree one, a candidate level set is a subset $S\subset\Phi^+$ on which a
linear functional can satisfy $f(\alpha)=1$ for every $\alpha\in S$.  Every
such $S$ contains a linearly independent subset with the same affine span.
This gives exhaustive small seeds and avoids scanning all $2^{24}$ subsets.
We keep the resulting list of $4862$ fine pairs.

\begin{tcolorbox}
\begin{algorithmic}[1]
\Require $\Phi^+$, the degree bound $d=1$, and the reviewed pair list $P$.
\If{$P$ is absent}
  \State Generate the independent affine-consistent seeds.
  \State Expand them by BFS, stopping at affine inconsistency or the degree
  bound, and write the resulting canonical list $P$.
\EndIf
\For{$(\Psi_0,\Psi_2)\in P$}
  \State Recompute the difference lattice, $\Psi_0$, and the closure axioms.
  \State Regard the pair as a terminal BFS state; do not generate children.
\EndFor
\For{each independent subset of the required degree}
  \State Check that its fine pair occurs in $P$.
\EndFor
\State Return $P$.
\end{algorithmic}
\end{tcolorbox}
\renewcommand{\figurename}{Algorithm}
\captionof{figure}{Fine-pair certification and fallback BFS}
\label{alg:F4-bounded-bfs}

There are $4862$ distinct fine pairs.  All $3002$ relevant independent subsets
occur in this list.  The cutoff is sound because adjoining a root to $\Psi_2$
strictly increases minimal degree.

\subsection{Exhaustive classification}

We apply the cone calculation to the $105$ possibilities for $K$ and the
$4862$ fine pairs.  Exactly four cases have $\varepsilon=0$; these are
Configurations I--IV.  All integral rank witnesses remain nonzero for
$p>12$.

For the four displayed grouped cone matrices the tuples
$(\#\mathrm{variables},\#\mathrm{relations},\operatorname{rank},
\mathrm{dimension})$ are
\[
(28,17,17,11),\quad(29,17,17,12),\quad
(30,14,14,16),\quad(32,8,7,25).
\]
Relations with the same target are one summed equation, never separate ideal
generators.  The implementation details and reviewed data are available at
\url{https://github.com/mocham/AlgEG/tree/main/code/v4}.
